\begin{frame}{그림으로 보는 ``추가 요소'' 하나}
  \begin{columns}[T]
    \begin{column}{0.48\textwidth}
      \centering
      \includegraphics[width=\linewidth]{ch02_3_bandit_vs_mdp.png}
      \vskip0.3em
      {\footnotesize 왼쪽: 밴딧(상태 1개, 자기 전이만).
      오른쪽: MDP(상태 2개, 실선으로 다른 상태로 전이).}
    \end{column}
    \begin{column}{0.5\textwidth}
      \begin{itemize}
        \item \textbf{왼쪽(밴딧)}: 상태 $s^*$ 하나 위에 세 개의
          \textbf{자기 전이}(팔 A, B, C)만 --- 어디로도
          ``넘어가지'' 않으므로, 지금 선택이 미래 선택지에
          영향을 줄 \textbf{통로가 전혀 없음}.
        \item \textbf{오른쪽(MDP)}: 상태 두 개($s_0, s_1$) ---
          $s_0$ 에서 한 행동이 $s_1$ 로 전이하고(실선),
          $s_1$ 에서 다시 $s_0$ 로 돌아옴(점선).
        \item \textbf{이 ``다른 상태로 넘어가는 실선''이 밴딧에는
          없는 유일한 추가 요소} --- Week 3의 $P(s'|s,a)$.
      \end{itemize}
    \end{column}
  \end{columns}
\end{frame}
